\section{Zentrale Begriffe und Frameworks}
\textless Text\textgreater

Risikomanagement ist teil aller Organisationsprozesse. Es gibt demnach eine stetige Evaluation der Risiken.

Im Vergleich zu klassischen Software- oder IT-Projekten ist das Risikomanagement in Data-Science-Projekten mit einer deutlich höheren Komplexität und Unsicherheit verbunden. Vor diesem Hintergrund ist eine klare Abgrenzung der relevanten Fachbegriffe notwendig, um die Besonderheiten dieses Projektkontexts sichtbar zu machen. (Holtkemper et al., 2023; Lahiri & Saltz, 2024 )

Während technologische Aspekte in Arbeitsgruppen häufig im Vordergrund stehen, wird unzureichendes Projektmanagement immer wieder als wesentliche Ursache für das Scheitern von Initiativen identifiziert. Ein effektives Risikomanagement als zentraler Bestandteil erfolgreichen Projektmanagements muss daher gezielt implementiert werden. Dazu gehört sowohl die Anwendung methodischer Rahmenwerke (Scrum oder CRISP-DM) als auch die Orientierung an spezifischen Standards (ISO, NIST, COSO etc.). So können typische Risiken in Data-Science-Projekten, wie etwa Overfitting oder systemischer Bias, frühzeitig erkannt und reduziert und gleichzeitig regulatorische Anforderungen der DSGVO oder des EU AI Acts, der als erste Verordnung seiner Art KI-Modelle drei Risikokategorien zuteilt, erfüllt werden. (Holtkemper et al., 2023; EU Artificial Intelligence Act, o.D.)

Die folgende Tabelle fasst die zentralen Begriffe zusammen, die für ein erfolgreiches Risikomanagement in Data-Science-Projekten besonders relevant sind.

\begin{center}

\begin{tabular}{|p{3cm}|p{4.9cm}|p{3cm}|p{4.9cm}|}
    \hline
    \textbf{Begriffe} & \textbf{Definitionen} & \textbf{Framework} & \textbf{Erläuterung} \\
    \hline
    Data Governance & Rahmenwerk, um qualitativ hochwertige Daten und die regelkonforme Nutzung sowie Verwahrung sicherstellen zu können & ISO & Internationale Organisation für Normungen \\
    \hline
    Overfitting & Überangepasstheit des Modells an Trainingsdaten, sodass mit neuen Datensätzen keine genauen Vorhersagen mehr getroffen werden können & ISO 31000 & Internationale Norm zur Identifikation und Vermeidung von Risiken \\
    \hline
    Bias & Vorurteile eines Modells aufgrund von homogenen oder veralteten Daten, was zu diskriminierenden Vorhersagen führen kann & ISO 21500 & Liefert Anleitungen für effektives Projektmanagement \\
    \hline
    Scope & Umfang eines Projekts in Bezug auf Aufwand, Ziele und Ergebnisse & ISO/IEC 27000 & Normreihe für Informationssicherheit, beinhaltet Basisnormen ISO 27001-27008 \\
    \hline
    Deployment & Bereitstellung eines entwickelten Modells & PMBOK & Project Management Body of Knowledge, Sammlung von Best Practices für Projektmanagement \\
    \hline
    Safeguards & Maßnahmenlisten, um Risiken vermeiden zu können & SCRUM & Methode des agilen Projektmanagements, gekennzeichnet durch eine iterative Arbeitsweise \\
    \hline
    Ease of Exploitation & Ausnutzbarkeit & COSO ERM -- Enterprise Risk Management Framework & Sammlung bewährter Methoden für Projekt- und Risikomanagement \\
    \hline
    Protection Effect & Schutzwirkung der Maßnahme & NIST Risk Management Framework & US-Standard für effektives Management von Sicherheitsrisiken \\
    \hline
    Information Asset & Informationswert & & \\
    \hline
\end{tabular}

\end{center}
